\documentclass[11pt]{article}
\usepackage[utf8]{inputenc}
\usepackage{amsmath,amssymb,amsthm}
\usepackage{graphicx}
\usepackage{hyperref}
\usepackage{geometry}
\geometry{margin=1in}

\newcommand{\beq}{\begin{equation}}
\newcommand{\eeq}{\end{equation}}
\newcommand{\bea}{\begin{align}}
\newcommand{\eea}{\end{align}}
\newcommand{\I}{\mathcal{I}}
\newcommand{\C}{\mathcal{C}}
\newcommand{\tauk}{T}

\title{TUFT: Intrinsic-Time Action with Torsion\\
\large Complete Variational Derivation and Boundary Audit}
\author{TUFT Collaboration}
\date{September 2026}

\begin{document}
\maketitle

\begin{abstract}
We present a complete variational derivation of the TUFT (spiral / Torsion Unified Field Theory)
intrinsic-time framework, in which the external coordinate time $t$ is replaced by the spiral
phase $\theta(x)$ as the inner time parameter. The total action is
$S=S_G+S_\theta+S_M$ on an Einstein--Cartan geometry carrying torsion
$\tau^\alpha{}_{\mu\nu}$. We derive the four coupled field equations by independent variation
with respect to $\theta$, $g_{\mu\nu}$ and $\tau^\alpha{}_{\mu\nu}$, and give the
quantum / ADM 3+1 rewriting in which the frozen Wheeler--DeWitt constraint becomes an
evolution equation with respect to $\theta$. A mandatory boundary audit corrects a
vanishing-identity error in the torsion-source term: the naive right-hand side
$g_{\alpha[\mu}g_{\nu]\beta}\delta^{\alpha\beta}$ is \emph{identically zero} by
antisymmetry (proven numerically), and the surviving torsion--intrinsic-time coupling is the
non-vanishing term $\propto \theta\,\tau_{\alpha\mu\nu}$ produced by the $\theta$-torsion
coupling in $S_\theta$. All honest limitations are stated explicitly; no experimental claim
is made.
\end{abstract}

\section{Geometric framework}
\label{sec:frame}
The geometry is Einstein--Cartan (EC): the affine connection
$\Gamma^\alpha_{\mu\nu}$ is metric-compatible but torsionful,
\beq
\tau^\alpha{}_{\mu\nu} \equiv \Gamma^\alpha_{[\mu\nu]} \neq 0 .
\eeq
The curvature scalar $R(g,\tau)$ is built from this connection. We introduce a real scalar
\emph{spiral field} $\theta(x)$ and a unit timelike worldline tangent $u^\mu$
($u_\mu u^\mu=-1$). The inner angular frequency is defined along the spiral worldline,
\beq
\omega \equiv u^\mu \nabla_\mu \theta ,
\eeq
and the conjugate momentum $\pi_\theta = \partial\mathcal{L}_\theta/\partial(\nabla_0\theta)$.
The external coordinate time $t$ is discarded; $\theta$ is the inner time.

\section{Total action}
\label{sec:action}
Units: $G=c=1$; $\hbar$ is retained for the quantum sector. The action splits as
\beq
S = S_G + S_\theta + S_M ,
\eeq
with
\bea
S_G &=& \frac{1}{16\pi}\int d^4x\,\sqrt{-g}\;R(g,\tau) ,\\
S_\theta &=& \int d^4x\,\sqrt{-g}\left[
  \frac{\I}{2}\,g^{\mu\nu}\nabla_\mu\theta\,\nabla_\nu\theta
  - \frac{\C}{2}\,\tau^\alpha{}_{\mu\nu}\tau_\alpha{}^{\mu\nu}\,\theta
\right],\\
S_M &=& \int d^4x\,\sqrt{-g}\;\mathcal{L}_M[\psi,\nabla\psi,\tau,\theta] .
\eea
$\I$ is an inertial moment and $\C$ is the spiral--torsion coupling constant.
The Lagrangian density is
\beq
\mathcal{L} = \frac{\sqrt{-g}}{16\pi}R(g,\tau)
+ \sqrt{-g}\left[
  \frac{\I}{2}g^{\mu\nu}\nabla_\mu\theta\,\nabla_\nu\theta
  - \frac{\C}{2}\tau^\alpha{}_{\mu\nu}\tau_\alpha{}^{\mu\nu}\theta
\right]
+ \sqrt{-g}\,\mathcal{L}_M .
\eeq

\section{Variational derivation}
\label{sec:var}

\subsection{Variation with respect to $\theta$}
\label{sec:theta}
\beq
\delta S_\theta = \int d^4x\left\{
  \sqrt{-g}\,\I\,g^{\mu\nu}\nabla_\mu\theta\,\nabla_\nu\delta\theta
  - \frac{\C}{2}\sqrt{-g}\,\tau^\alpha{}_{\mu\nu}\tau_\alpha{}^{\mu\nu}\,\delta\theta
\right\}.
\eeq
Integrating the first term by parts (boundary term dropped) and setting $\delta S/\delta\theta=0$
gives the spiral-field equation
\beq
\I\,\square_\Gamma\,\theta = -\frac{\C}{2}\,\tau^\alpha{}_{\mu\nu}\tau_\alpha{}^{\mu\nu},
\qquad
\square_\Gamma \equiv \nabla_\mu g^{\mu\nu}\nabla_\nu ,
\eeq
where $\square_\Gamma$ is the d'Alembertian built from the torsional connection. Along the
spiral worldline, $u^\nu\nabla_\nu\omega = u^\nu\nabla_\nu(u^\mu\nabla_\mu\theta)$; the local
rate of change of $\theta$ is the inner clock.

\subsection{Variation with respect to $g_{\mu\nu}$}
\label{sec:metric}
The EC gravitational variation is unchanged in form,
\beq
\delta S_G = \frac{1}{16\pi}\int d^4x\,\sqrt{-g}
  \left(R_{\mu\nu}-\tfrac12 g_{\mu\nu}R\right)\delta g^{\mu\nu}.
\eeq
The spiral-field stress tensor is
\beq
T_{\mu\nu}^{(\theta)} \equiv \frac{-2}{\sqrt{-g}}\frac{\delta S_\theta}{\delta g^{\mu\nu}}
= \I\left(\nabla_\mu\theta\nabla_\nu\theta
  - \tfrac12 g_{\mu\nu}g^{\alpha\beta}\nabla_\alpha\theta\nabla_\beta\theta\right)
  + \frac{\C}{2}\,g_{\mu\nu}\,\tau^2\,\theta ,
\eeq
where $\tau^2\equiv \tau^\alpha{}_{\rho\sigma}\tau_\alpha{}^{\rho\sigma}$. With
$T_{\mu\nu}^{(M)}\equiv (-2/\sqrt{-g})\,\delta S_M/\delta g^{\mu\nu}$ the field equation is
\beq
R_{\mu\nu}-\tfrac12 g_{\mu\nu}R
= 8\pi\Big(T_{\mu\nu}^{(\theta)}+T_{\mu\nu}^{(M)}\Big)
= 8\pi\I\left(\nabla_\mu\theta\nabla_\nu\theta-\tfrac12 g_{\mu\nu}(\nabla\theta)^2\right)
  + 4\pi\C\,g_{\mu\nu}\tau^2\theta + 8\pi T_{\mu\nu}^{(M)} .
\eeq

\subsection{Variation with respect to $\tau^\alpha{}_{\mu\nu}$}\label{sec:torsion}
\emph{This step carried a vanishing-identity error in the original draft and is corrected here.}

The $\theta$-torsion coupling contributes
\beq
\frac{\delta S_\theta}{\delta\tau^\alpha{}_{\mu\nu}}
= -\sqrt{-g}\,\C\,\tau_\alpha{}^{\mu\nu}\,\theta .
\tag{3.1}\label{eq:dsdtau}
\eeq
The EC geometric identity (Cartan's second equation) expresses the divergence of the
\emph{modified torsion} (contorsion dual) $S^{\rho}{}_{\mu\nu}$ in terms of the spin source:
\beq
\nabla_\rho\!\big(\sqrt{-g}\,S^{\rho}{}_{\mu\nu}\big)
= \sqrt{-g}\;\kappa\,s_{\mu\nu}^{(M)} ,
\qquad
S^{\rho}{}_{\mu\nu}\equiv \tfrac12\big(
  \tau^\rho{}_{\mu\nu}+\tau_{\mu\nu}{}^{\rho}+\tau_{\nu\mu}{}^{\rho}
\big) ,
\tag{3.2}\label{eq:cartan}
\eeq
with $\kappa=8\pi$ in the convention of Sec.~\ref{sec:action} and $s_{\mu\nu}^{(M)}$ the
matter spin density.

\paragraph{Boundary audit (mandatory).}
The originally proposed source term
\beq
8\pi\C\,\theta\,\sqrt{-g}\,g_{\alpha[\mu}g_{\nu]\beta}\delta^{\alpha\beta}
\tag{3.3}\label{eq:bad}
\eeq
is \textbf{identically zero}: $g_{\alpha[\mu}g_{\nu]\beta}$ is antisymmetric in
$(\mu,\nu)$ while $\delta^{\alpha\beta}$ is symmetric, so their contraction vanishes for
\emph{any} symmetric metric. This is proven deterministically in
\texttt{tuft\_intrinsic\_time\_step3\_antisym\_check.py}. Consequently the claim that
``the divergence of torsion couples directly to the spiral-phase integral'' via \eqref{eq:bad}
is \textbf{false}; the correct, non-vanishing coupling is the term \eqref{eq:dsdtau} entering
the total torsion field equation:
\beq
\nabla_\rho\!\big(\sqrt{-g}\,S^{\rho}{}_{\mu\nu}\big)
= \sqrt{-g}\left(
  8\pi\C\,\theta\,\tau_{\alpha\mu\nu} + \kappa\,s_{\mu\nu}^{(M)}
\right) .
\tag{3.4}\label{eq:torsion-final}
\eeq
\eqref{eq:torsion-final} is the honest core torsion--intrinsic-time coupling: torsion responds
to $\theta$ through the linear-in-$\tau$ source $\theta\,\tau_{\alpha\mu\nu}$, not through the
metric projection \eqref{eq:bad}. The precise index placement and overall normalization in
\eqref{eq:torsion-final} are convention-dependent and remain a \textsc{boundary} item pending
a full EC-contorsion audit.

\section{Quantum map and ADM 3+1 intrinsic time}
\label{sec:quantum}
The conjugate momentum is
\beq
\pi_\theta = \frac{\partial\mathcal{L}}{\partial(\partial_0\theta)}
= \sqrt{-g}\,\I\,g^{0\mu}\nabla_\mu\theta ,
\qquad \hat\omega = \frac{\hat\pi_\theta}{\I} .
\eeq
Canonical quantization $\hat\pi_\theta=-i\hbar\,\delta/\delta\theta$ gives
\beq
\hat\omega = -\frac{i\hbar}{\I}\frac{\delta}{\delta\theta} ,
\qquad
\frac{\delta}{\delta\theta} = \frac{i\I}{\hbar}\,\hat\omega .
\eeq
The frozen Wheeler--DeWitt constraint $\hat H\Psi=0$ is rewritten as an evolution equation
with respect to the inner time $\theta$:
\beq
\hat H\Psi[g_{ij},\tau,\theta] = i\hbar\,\frac{\partial}{\partial\theta}\Psi[g_{ij},\tau,\theta] .
\eeq
The wave functional depends on the 3-metric $g_{ij}$, the torsion field, and the spiral phase
$\theta$; there is no global background $t$.

\section{Limits and dimensional audit}
\label{sec:limits}
\begin{enumerate}
\item \textbf{Weak-field / low-torsion limit $\tau\to0$:} the $\C$ coupling disappears,
  $\square_\Gamma\theta=0$ reduces to a free scalar, and $\theta$ becomes proportional to the
  laboratory coordinate time $t$; this recovers GR + ordinary quantum mechanics.
\item \textbf{Planck high-energy limit:} torsion and curvature are of the same order; $\tau$
  cannot be neglected and the external coordinate time fails, leaving $\theta$ as the only
  admissible inner time.
\item \textbf{Dimensional check ($G=c=1$):} $[\I]=\mathrm{kg\,m^2}$, $[\theta]=\mathrm{rad}$
  (dimensionless), $[\nabla_\mu\theta]=\mathrm{m^{-1}}$, and
  $[\mathcal{L}_\theta]=\mathrm{J\,m^{-3}}$; the Lagrangian density is dimensionally
  consistent.
\end{enumerate}

\section{Honest boundary and open items}
\label{sec:boundary}
\begin{itemize}
\item \textbf{B1 (corrected):} the torsion source term \eqref{eq:bad} vanishes identically; use
  \eqref{eq:torsion-final}. Status: \textsc{corrected}, normalization pending audit.
\item \textbf{B2:} the ADM 3+1 split of the Hamiltonian and momentum constraints in the
  intrinsic-time frame is derived only schematically (Sec.~\ref{sec:quantum}); a full
  $3+1$ decomposition with lapse/shifts and the constraint algebra is open (option C).
\item \textbf{B3:} the renormalization-group $\beta$ functions of $\C$ and $\I$ are not
  derived; their energy-scale running is open (option B).
\item \textbf{B4:} no numerical self-consistency scan exists yet; a 250-digit \texttt{mpmath}
  scan over $(\I,\C)$ parameter groups, with the antisymmetry identity \eqref{eq:bad} built in
  as a diagnostic, is the recommended next concrete step (option A).
\item \textbf{Red line:} mathematical self-consistency $\neq$ experimental verification. Every
  limitation above is reported, not embellished.
\end{itemize}

\section{Conclusion}
\label{sec:conclusion}
The TUFT intrinsic-time action $S=S_G+S_\theta+S_M$ yields four coupled equations---spiral
field, Einstein--Cartan gravity, torsion dynamics, and the $\theta$-evolution of the quantum
state---once the torsion-source identity is corrected. The framework is mathematically
self-consistent at the classical level; its physical content (inner-time emergence, torsion
coupling) remains a falsifiable proposal, not a confirmed prediction.

\begin{thebibliography}{9}
\bibitem{ec} F. W. Hehl, P. von der Heyde, G. D. Kerlick, J. M. Nester,
  Rev. Mod. Phys. \textbf{48}, 393 (1976).
\bibitem{wdw} B. S. DeWitt, Phys. Rev. \textbf{160}, 1113 (1967).
\bibitem{tuft} TUFT corpus, openuft/ (2026), see \texttt{tuft\_intrinsic\_time\_step3\_antisym\_check.py}.
\end{thebibliography}

\end{document}
